% Irodov Part III §3.3 131--140; source PDF pp.119--120.
% The PDF printed page 110 repeats the title of problem 130 as an extra "131";
% the second "131" is the capacitor problem matched by the original answer 131.
\electrocapproblem{131}{110}
电容为 $C_1=1.0\,\mathrm{\mu F}$、
初始电压 $V=300\,\mathrm V$ 的电容器，
与一只未充电的电容器
$C_2=2.0\,\mathrm{\mu F}$ 并联相接。
求系统达到平衡后的电场能增量，并解释能量去向。
\exsolution
连接前总电荷
$Q_0=C_1V$，总电场能
$W_0=\frac12C_1V^2$。
并联后两只电容器电势相等且总电荷守恒：
\[
V_f=\frac{Q_0}{C_1+C_2}=\frac{C_1V}{C_1+C_2}.
\]
末态电场能
$W_f=\frac12(C_1+C_2)V_f^2
=\frac12 C_1^2V^2/(C_1+C_2)$，
所以
\[
\boxed{\Delta W=W_f-W_0
=-\frac12V^2\frac{C_1C_2}{C_1+C_2}}.
\]
数值：
\[
\boxed{\Delta W=-0.030\,\mathrm J=-30\,\mathrm{mJ}}.
\]
消失的电场能在实际连接导线中
变为焦耳热、电磁辐射等，
不会凭空消失。
\excheck 原版 PDF 的正文将第130题重复印成额外
“131”，此题是第二次出现的“131”且与答案131对应。
书后数值若显示“$0.03\,\mathrm{mJ}$”
则与 $1\mu\mathrm F$、300V 的计算相差1000倍；
本稿保留正确 SI 算值并标勘误。

\electrocapproblem{132}{110}
图3.30：两只电容 $C$、
中间电容 $C_0$ 与电池 $\mathcal E$
构成电路。开关 Sw 从位置1拨到2。
求开关过程中产生的焦耳热。
\exsolution
把开关前后相应节点看作
两个通过电容器耦合的等势导体，
对于固定电源电动势，
电容电路的\emph{电势跳变}激发一次电荷再分配。
图3.30左右两个相等电容 $C$
从公共上导体接到由电池保持电压差
$\mathcal E$ 的下方两节点。
设开关连接的中间电容 $C_0$
下端初连左节点、末连右节点。
切换引起的中间支路电位跃变为 $\mathcal E$；
消除其公共模式后，
参与再分配的等效电容
\[
C_{\rm tr}=\frac{2CC_0}{2C+C_0}.
\]
一次从静态到静态的切换
消耗的电场—电源可用能
$\frac12C_{\rm tr}(\Delta V)^2$，
于是
\[
\boxed{Q=\frac{CC_0}{2C+C_0}\,\mathcal E^2}.
\]
\excheck 本题求的是有微小实际电阻时最终释放的
总焦耳热，理想零电阻模型只提供初末静态电荷；
若开关瞬时切换，过程可含瞬态电磁辐射。

\electrocapproblem{133}{110}
图3.31：一只电容 $C$ 通过换向开关
先接与电动势 $\mathcal E_1$ 对应的节点，
再从位置1切至位置2，
后者改变电容两端的电位差，
求电路释放的热量。
\exsolution
设切换前电容器电压为 $U_i$，
切换后电压为 $U_f$。
对于理想电源维持的最终电压状态，
通过微弱电阻缓慢耗散的电荷再分配过程，
其热量满足普适恒等式
\[
Q=\frac12C(U_f-U_i)^2.
\]
由原书图3.31两个源的串联位置及开关端点，
开关从1切到2时电容器电压\emph{净改变}
$U_f-U_i=\mathcal E_2$，
所以
\[
\boxed{Q=\frac12C\mathcal E_2^2}.
\]
\excheck 结果与 $\mathcal E_1$ 无关，
因为焦耳热只由切换产生的电压差平方决定；
本结论使用图中电源的具体极性与接线。

\electrocapproblem{134}{110}
两层同心薄金属球壳半径
$R_1<R_2$，电荷分别为 $q_1,q_2$。
求各球壳静电自能 $W_1,W_2$，
两球壳相互作用能 $W_{12}$，
以及总电场能。
\exsolution
单个薄球壳电场能为
$W_{\rm self}=q^2/(8\pi\varepsilon_0R)$，
所以
\[
\boxed{W_1=\frac{q_1^2}{8\pi\varepsilon_0R_1}},
\qquad
\boxed{W_2=\frac{q_2^2}{8\pi\varepsilon_0R_2}}.
\]
外球壳在内球壳位置的电势恒为
$kq_2/R_2$，故交互能
\[
\boxed{W_{12}=\frac{q_1q_2}{4\pi\varepsilon_0R_2}}.
\]
相加：
\[
\boxed{W=
\frac1{8\pi\varepsilon_0}
\left(\frac{q_1^2}{R_1}+
\frac{q_2^2+2q_1q_2}{R_2}\right)}.
\]
\excheck 交互能只计入一次，
不应在 $W_1,W_2$ 内重复算交互项。

\electrocapproblem{135}{110}
半径为 $R$ 的实心绝缘球，
体电荷均匀且总电荷为 $q$，
周围真空。求总静电自能
及球内储能 $W_{\rm in}$ 与球外储能
$W_{\rm out}$ 的比值。
\exsolution
高斯定律给球内
$E(r)=qr/(4\pi\varepsilon_0R^3)$，
球外
$E(r)=q/(4\pi\varepsilon_0r^2)$。
场能密度
$w=\varepsilon_0E^2/2$。
球内积分：
\[
W_{\rm in}=\int_0^R
\frac{\varepsilon_0}{2}
\frac{q^2r^2}{16\pi^2\varepsilon_0^2R^6}
4\pi r^2\,\mathrm dr
=\frac{q^2}{40\pi\varepsilon_0R}.
\]
球外积分：
\[
W_{\rm out}=\int_R^\infty
\frac{\varepsilon_0}{2}
\frac{q^2}{16\pi^2\varepsilon_0^2r^4}
4\pi r^2\,\mathrm dr
=\frac{q^2}{8\pi\varepsilon_0R}.
\]
于是
\[
\boxed{W=W_{\rm in}+W_{\rm out}
=\frac{3q^2}{20\pi\varepsilon_0R}},
\qquad
\boxed{\frac{W_{\rm in}}{W_{\rm out}}=\frac15}.
\]
\excheck 球体内能只占总场能 $1/6$；
若把实心球按金属球壳处理将漏掉球内能。

\electrocapproblem{136}{110}
点电荷 $q=3.0\,\mathrm{\mu C}$
位于球心，
$R_1=a=250\,\mathrm{mm}$ 至
$R_2=b=500\,\mathrm{mm}$ 的球壳空间
填满相对介电常数 $\varepsilon=3.0$
的均匀介质。
求仅存储在这层\emph{介质体积}内的静电场能。
\exsolution
介质中的径向位移
$D=q/(4\pi r^2)$，
电场
$E=q/(4\pi\varepsilon_0\varepsilon r^2)$。
场能密度
$w=\frac12\bm E\cdot\bm D$，
故
\[
W=\int_a^b \frac{q^2}
{32\pi^2\varepsilon_0\varepsilon r^4}
4\pi r^2\,\mathrm dr
=\frac{q^2}{8\pi\varepsilon_0\varepsilon}
\left(\frac1a-\frac1b\right).
\]
数值
\[
\boxed{W\simeq2.7\times10^{-2}\,\mathrm J
=27\,\mathrm{mJ}}.
\]
\excheck 本题只求介质\emph{球壳内}场能，
不包括 $r<a$ 与 $r>b$ 空间的能量。

\electrocapproblem{137}{110}
均匀带电薄球壳总电荷固定为 $q$，
半径从 $R_1$ 缓慢扩大为 $R_2>R_1$。
求静电斥力做功。
\exsolution
均匀带电薄球壳自能为
$W(R)=q^2/(8\pi\varepsilon_0R)$。
自由电荷守恒，静电力做的功
$A=W(R_1)-W(R_2)$，故
\[
\boxed{A=\frac{q^2}{8\pi\varepsilon_0}
\left(\frac1{R_1}-\frac1{R_2}\right)}.
\]
\excheck 扩张时场能下降，电力做正功；
这里忽略机械表面张力和其它非电作用。

\electrocapproblem{138}{111}
总电荷固定为 $q$ 的均匀薄球壳
半径从 $R_1$ 增至 $R_2$，
球心另外放有点电荷 $q_0$。
求电场力对球壳扩张所做功。
\exsolution
壳自身能量
$W_{\rm shell}=q^2/(8\pi\varepsilon_0R)$。
球心点电荷与壳的相互作用能为
$W_{\rm int}=q_0\varphi_{\rm shell}(0)
=q_0q/(4\pi\varepsilon_0R)$。
点电荷自能独立于 $R$，在能量差中相消。
因此
\[
A=[W_{\rm shell}(R_1)+W_{\rm int}(R_1)]
-[W_{\rm shell}(R_2)+W_{\rm int}(R_2)].
\]
整理：
\[
\boxed{A=
\frac{q(q+2q_0)}{8\pi\varepsilon_0}
\left(\frac1{R_1}-\frac1{R_2}\right)}.
\]
\excheck 若 $q_0$ 与 $q$ 异号且足够大，
作用力可能反而阻碍扩张，此时 $A<0$。

\electrocapproblem{139}{111}
均匀带电薄球壳的面电荷密度为 $\sigma$。
利用能量守恒求其表面每单位面积所受
向外静电压力（法向力模）。
\exsolution
半径为 $R$ 的球壳总电荷
$q=4\pi R^2\sigma$。
以总电荷 $q$ 固定为前提，
球壳自能
$W=q^2/(8\pi\varepsilon_0R)$。
若半径增大 $\mathrm dR$，
体积增量
$\mathrm dV=4\pi R^2\mathrm dR$。
向外压力 $p_e$ 做功
$p_e\mathrm dV=-\mathrm dW$，所以
\[
p_e=-\frac1{4\pi R^2}\frac{\mathrm dW}{\mathrm dR}
=\frac{q^2}{32\pi^2\varepsilon_0R^4}.
\]
利用
$\sigma=q/(4\pi R^2)$：
\[
\boxed{p_e=\frac{\sigma^2}{2\varepsilon_0}}.
\]
\excheck 变形求导时应保持总电荷 $q$
不变，不能误保持面电荷密度 $\sigma$ 不变。

\electrocapproblem{140}{111}
球心 O 放有点电荷 $q$，
其外有一层整体不带净电荷的球形导体壳，
内外半径为 $a<b$，并有小孔（图3.32）。
求通过小孔把电荷从中心缓慢移至无穷远
需要的外功。
\exsolution
以导体壳的感应电荷解为基础：
中心点电荷使内球面感应 $-q$，
外球面感应 $+q$。
初态电场在 $0<r<a$ 和 $r>b$
分别存在，中间导体内部为零。
若将点电荷移至无穷远，
球壳总电荷仍为零，极限态无电场。
两态与点电荷自身发散自能无关的
\emph{可变有限场能}差，
可由空心球壳在两个半径之间的
缺失场能计算：
\[
A=\frac12\int_a^b
\frac{q^2}{(4\pi\varepsilon_0r^2)^2}
\varepsilon_0(4\pi r^2)\,\mathrm dr.
\]
故
\[
\boxed{A=\frac{q^2}{8\pi\varepsilon_0}
\left(\frac1a-\frac1b\right)}.
\]
\excheck 假设小孔足够小，不改变球对称场近似；
真实经过孔口时局部场会畸变，
但端点能量差在该理想极限不变。
